%%%%

%%%   Самарский государственный технический университет

%%%   Кафедра прикладной математики и информатики

%%%

%%%   Преамбула для статьи в журнал "Вестник Самарского государственного

%%%   технического университета. Серия: «Физико-математические науки»"

%%%   Ориентирована под MikTEx 2.4 for Win32

%%%

%%%   Хотя сам журнал собирается под GNU/Linux на дистрибутиве TeXLive



\documentclass[11pt,a4paper,twoside]{article}



\usepackage[cp1251]{inputenc}

\usepackage[T2A]{fontenc}

\usepackage[english,russian]{babel}

\usepackage{samgtubib}

\usepackage[dvips]{graphicx}

\usepackage[multiple]{footmisc}



\usepackage{samgtu}

\renewcommand{\VestnikYear}{2009} % Год издания Вестника

\renewcommand{\VestnikNumber}{2\,(19)} % Номер Вестника

\renewcommand{\VestnikMonth}{Ноябрь} % Месяц выхода Вестника

\renewcommand{\VestnikData}{01.11.09} % Дата подписания Вестника в печать

\renewcommand{\VestnikUPL}{26,64} % Усл. печ. л.

\renewcommand{\VestnikUIL}{25,73} % Уч.-изд. л.

\renewcommand{\VestnikRegNumber}{479/09} % Регистрационный номер

\renewcommand{\VestnikOrder}{994} % Номер заказа







%

%\usepackage{pst-barcode}

%\usepackage{auto-pst-pdf}

%%

%%\newcommand{\totop}[1]{\raisebox{6pt}[1pt][2pt]{#1}} 

%%\newcommand{\tottop}[1]{\raisebox{12pt}[1pt][2pt]{#1}} 

%%\newcommand{\totttop}[1]{\raisebox{18pt}[1pt][2pt]{#1}} 

%%\newcommand{\tottttop}[1]{\raisebox{24pt}[1pt][2pt]{#1}} 

%%\newcommand{\totttttop}[1]{\raisebox{30pt}[1pt][2pt]{#1}} 

%%\newcommand{\tobot}[1]{\raisebox{-6pt}[1pt][2pt]{#1}} 

%%\newcommand{\tobbot}[1]{\raisebox{-12pt}[1pt][2pt]{#1}} 

%%\newcommand{\tobbbot}[1]{\raisebox{-18pt}[1pt][2pt]{#1}}

%

%\graphicspath{{./00PIC/}}

%

%%\samgtupubl % для генерации pdf, отдаваемого в типографию СамГТУ

%\pdfcrop % обрезка pdf для размещения на math-net.ru

%\draft % На корректуре

%\filllastpage % Последняя страница не заполнена не полностью

%



\vsgtu{1437}



\FirstPage{785}

\LastPage{798}

%%

%%

%\begin{document}





%\hfill \begin{pspicture}(1in,1in)

%\psbarcode{http://dx.doi.org/10.14498/vsgtu1437}{}{qrcode}

%\end{pspicture}

%\vspace{-1in}



% Номер УДК

\udk{004.434; 004.453}%Программное управление вычислительными машинами. Программы обслуживания

\msc{68M14, 68M20}% Коды MSC см. http://www.ams.org/msc/



\rutitle{Инструментальное программное обеспечение для~разработки и поддержки исполнения приложений научных вычислений в~кластерных~системах}% Полный заголовок

{Инструментальное программное обеспечение\ldots}% Заголовок, который идёт в верхний колонтитул



\ruauthor{Ю.~С.~Артамонов, С.~В.~Востокин}{А\,р\,т\,а\,м\,о\,н\,о\,в~Ю.~С.,\ В\,о\,с\,т\,о\,к\,и\,н~С.~В.}



\ruaffil{\rusgau.}



%\email{E-mails}{artamonov@about.me, easts@mail.ru}



\ruauthordetails{2}{\fullauthor{\rupid{112159}{Юрий Сергеевич Артамонов}} \regalia{\mem{artamonov@about.me}} \position{аспирант}

\dept{каф. информационных систем и технологий}\\

\fullauthor{\rupid{27041}{Сергей Владимирович Востокин}}

\regalia{д.т.н., проф.; \mem{easts@mail.ru}; автор, ведущий переписку} 

\position{профессор} \dept{каф. информационных систем и технологий} 

}



\enauthordetails{2}{\fullauthor{\enpid{112159}{Yuriy S. Artamonov}}  \regalia{\mem{artamonov@about.me}} \position{Postgraduate Student}\dept{Dept. of Information Systems \& Technology} \\

\fullauthor{\enpid{27041}{Sergey V. Vostokin}}\regalia{Dr. Techn. Sci.;  \mem{easts@mail.ru}; Corresponding Author} \position{Professor} \dept{Dept. of Information Systems \& Technology}

}





\rukeywords{инструмент, сервис, библиотека, параллельное программирование, окружение, запуск, мониторинг, вычисления, кластер}



\ruabstract{\textbf{Мотивация:} Для разработки приложений научных вычислений существует множество различных инструментов. Большинство из них ориентированы на сам процесс написания программ, но часто требуются приложения для организации процесса вычислений и поддержки командной разработки. Описана специфика разработки приложений научной направленности, сделан акцент на характерных проблемах разработки такого ПО. \textbf{Классификация систем управления вычислительными задачами:} Приводится классификация систем по способу организации вычислений и уровню абстракции вычислений от физического оборудования.  \textbf{Инструменты разработки Templet:} Рассматривается инструментарий для разработки приложений, включающий в себя библиотеки параллельных вычислений, сервис запуска и отслеживания задач, подсистему мониторинга состояния кластера. Тесное взаимодействие инструментов позволяет эффективно организовать работу команды над приложением научной направленности. \textbf{Решение прикладных задач при помощи инструментов Templet:} Инструментарий применяется для решения практических задач в области моделирования поведения многомерных динамических систем. Показан подход, позволяющий разделить работу над приложением на системный и прикладной уровни. \textbf{Заключение:} Сделан вывод о возможностях применения техник проектирования и преимуществах, которые даёт использование инструментария.}



\entitle{Tooling software for development and execution support of scientific computing applications in~cluster systems}



\enauthor{Yu.~S.~Artamonov, S.~V.~Vostokin}{A\,r\,t\,a\,m\,o\,n\,o\,v~Yu.~S.,\ V\,o\,s\,t\,o\,k\,i\,n~S.~V.}

%

\enaffil{\ensgau.}





\enabstract{\textbf{Rationale:} Many different tools exist for development of scientific computing applications. Most of them are focused on the process of writing software code, but often there is a need for applications that organize the computation process and support team development. The article describes application development specifics in the field of science-oriented computing and highlights individual issues in the development of such software. \textbf{Classification of task management systems:} The systems are classified by means of computing process organization and the layer of hardware abstraction. \textbf{Templet development tools:} The tools for application development considered in the article include parallel programming libraries, a task running and monitoring service and the monitoring subsystem for SSAU cluster. Close interaction between these tools enables effective teamwork for scientific application development. \textbf{Applied problems solved by Templet tools:} Tooling is used to solve practical issues in the field of modeling multi-dimensional dynamic systems behavior. The article demonstrates an approach that splits application development into system-level and applied development layers.  \textbf{Conclusion:} The article concludes about the use of design techniques and the benefits provided by software development tools.}



\enkeywords{tooling, service, library, parallel programming, environment, monitoring, computing, cluster}



\date{18/V/2015} % Дата поступления статьи в редакцию.

\revisiondate{27/VI/2015} % В окончательном виде

\accepteddate{08/VIII/2015}



%%%%%%%%%%%%%%%%%%%%%%%%%%%%%%%%%%%%%%%%%%%%%%%%%%%%%%%%%%%%%%%%%%%%%%%%%%%%%%%%%%%%



\makerutitle



%Каждый пункт статьи должен начинаться с команды Название пункта

%после названия точку ставить не нужно. Пункты автоматически нумеруются.

%Если пункт нумеровать не нужно, то следует использовать в команде

%\Section необязательный параметр [n]:

%\Section[n]{Имя пункта}

%Введение и заключение обычно не нумеруются.







\Section[n]{Мотивация}

Для разработки программ существует множество инструментов "--- среды разработки, компиляторы и трансляторы, библиотеки, парсеры, отладчики и многое другое. 

Практически все они ориентированы на создание конечного продукта и лишь немногие "--- на итеративную разработку и~поддержку. 

Отличительными особенностями приложений для научных вычислений являются, с одной стороны, короткий жизненный цикл приложения, при котором часто не требуется отдельная документация, тестирование в различных ОС, оформление графического интерфейса и другие финальные этапы разработки программного обеспечения (ПО), с~другой "--- большое количество вносимых изменений, требующее запуска приложения и получения результатов \cite[p.~50--54]{Keyes2013Challenges}. 

Программы научных вычислений очень специализированы и довольно сильно зависят от целевого окружения (кластера, сети, систем хранения данных), что делает разработку таких приложений непохожей на разработку системного и прикладного ПО.



Помимо специфики самих приложений стоит отметить ряд характерных проблем, возникающих в процессе работы над проектом:

\begin{enumerate}

\item растущая сложность алгоритмов и архитектуры приложений;

\item окружения, в которых производится запуск, зачастую имеют нестандартный набор ПО и периферийных устройств;

\item сложный и длительный процесс развёртывания ПО в целевом окружении;

\item для тестирования может потребоваться производительное окружение;

\item оборудование и инфраструктура могут быть уникальны и их нельзя воссоздать для тестирования;

\item необходимо соблюдать регламент работы с целевым окружением, определённое время работы на кластере и лимит нагрузки;

\item доступ к окружению может быть ограничен локальной сетью организации.

\end{enumerate}



Под окружением здесь понимается некоторый программно-аппаратный комплекс: совокупность физических или виртуальных ЭВМ, каналов связи, периферийных устройств и ПО, необходимого для работы системы. Например, это может быть кластер серверов, единичный сервер или же Desktop Grid-система.



Процесс организации работы над приложением также может быть сложным, поскольку одно приложение может разрабатываться целым коллективом. Причём командная разработка всё чаще требуется при работе над приложениями научных вычислений. Во многом это обусловлено тем, что команда может разрабатывать многие модули и подсистемы параллельно, к тому же в команде часто требуются как специалисты по предметной области, так и системные программисты, специалисты по параллельным вычислениям.



Отметим наиболее важные аспекты командной работы над приложениями:

\begin{enumerate}

\item[1.] Ни один крупный проект сейчас не может обойтись без использования систем контроля версий (VCS), ведь в процессе разработки код проекта претерпевает значительные изменения, их все требуется отслеживать. Обязательно должна быть возможность вернуться к более ранней версии \cite{Sullivan2009RCS}.

\item[2.] В проекте должен быть простой и эффективный механизм сборки приложения без использования графических утилит, поддерживающий \linebreak управление зависимостями и различные целевые архитектуры \cite{Neitsch2012Build}. 

Следует отметить, что для вычислений всё чаще используются языки, которым не требуется компиляция в машинно"=зависимый код: Java, Py\-thon, R, Lisp, Scheme. Использование этих языков существенно упрощает работу над проектом, но не стоит забывать о накладных расходах при исполнении в виртуальной машине (на работу интерпретатора).

\item[3.] Поскольку тестовый запуск может требовать значительных ресурсов, большое значение приобретают модульные и интеграционные тесты \cite{Ettl2012CI}. Причём такие тесты должны запускаться довольно часто и, по возможности, автоматически при поступлении новых изменений от участников проекта.

\item[4.] Крайне важно разделить тестовые запуски и запуски вычислений для получения конечных результатов. Следует организовать тестовое окружение, если это возможно, и тестовые запуски проводить в нём. Запуск приложения на суперкомпьютере для целей тестирования и отладки требует части ценных разделяемых ресурсов, что нерационально, кроме того, результаты тестирования могут быть получены с существенной задержкой в случае загруженности основных ресурсов.

\end{enumerate}



Для решения этих проблем разработано множество инструментов, позволяющих организовать работу над проектом, упростить запуск приложений и обработку результатов. В этой работе речь пойдёт о системах управления вычислительными задачами как основном инструменте ресурсоёмких научных вычислений.



\smallskip



\Section{Классификация систем управления вычислительными задачами}

Системы управления вычислительными задачами можно разделить на несколько специализированных классов:

\begin{enumerate}

\item низкоуровневые системы управления распределёнными ресурсами (пакетные системы);

\item Desktop Grid-системы;

\item системы виртуализации вычислительных ресурсов;

\item высокоуровневые платформы распределённых вычислений.

\end{enumerate}



\smallskip \textit{Низкоуровневые системы управления распределёнными ресурсами} призваны утилизировать доступные вычислительные ресурсы и максимизировать пропускную способность, управляя процессом назначения физических ресурсов пользовательским задачам. Пакетные системы управляют такими физическими ресурсами, как процессорное время, оперативная память, дисковое пространство и пропускная способность сети. К таким системам относятся \cite{Yan2008Comparative}: Platform LSF, Portable Batch System (PBS), Sun Grid Engine (SGE), IBM Load Leveler и Condor. Пакетные системы обычно имеют следующие подсистемы:

\begin{enumerate}

\item подсистема управления задачами;

\item подсистема управления физическими ресурсами;

\item подсистема планирования и управления очередями.

\end{enumerate}

Подсистема управления задачами "--- интерфейс для пользователей, которые взаимодействуют с системой, запрашивая ресурсы для задач и контролируя статус их исполнения. Подсистема управления физическими ресурсами решает две основных задачи: контроль использования ресурсов и сбор статистики их использования. В пакетных системах назначение ресурсов для задач зависит от политики планирования и использования ресурсов. Эта политика определяется администратором системы и является конфигурацией для подсистемы планирования и управления очередями. 

Основные достоинства пакетных систем:

\begin{enumerate}

\item пригодны для организации вычислений на различном оборудовании и операционных системах;

\item хорошо подходят для большого числа перезапусков задач;

\item позволяют организовать работу множества пользователей;

\item документированы, широко известны и повсеместно применяются.

\end{enumerate}

Недостатки:

\begin{enumerate}

\item требуется обучение пользователей, поскольку им необходимо знать особенности физического оборудования кластера для запуска задач, такие как число ядер процессоров, оперативная память, скорость чтения и записи данных, топология сети;

\item пакетные задания трудно отлаживать;

\item зависшие задачи могут держать занятыми ресурсы всё отведённое им время.

\end{enumerate}



\smallskip \textit{Desktop Grid-системы:} как и пакетные системы приложения desktop-grid, 

нацелены на утилизацию вычислительных ресурсов, однако они существенно отличаются от них. 

Пакетные системы задействуют оборудование организаций: суперкомпьютеры, кластеры и другое дорогостоящее оборудование. В~отличие от пакетных систем, Desktop Grid задействуют обычные компьютеры \cite{Anderson2004Boinc}, а связь между ними организуется через сеть Интернет, кроме того,  компьютеры могут быть выключены или отключены от сети. 

Участники проектов Desktop Grid не являются экспертами и разработчиками ПО, они участвуют в~вычислениях, поскольку им это интересно. Проекты не контролируют своих участников и не могут предотвратить злонамеренное поведение, например, предоставление ложных результатов. 



Системы Desktop Grid могут иметь клиентское приложение, реализованное при помощи веб-апплета или отдельного настольного приложения. Если клиентское приложение реализовано в виде веб-апплета, то участникам проекта достаточно открыть веб-страницу приложения и запустить исполнение заданий. Апплет реализован на языке Java и может использовать ресурсы настольного компьютера с согласия пользователя. К таким приложениям относятся: Charlotte, Javelin и Bayanihan.  В~случае использования отдельного настольного приложения пользователю потребуется установить и запустить клиентское приложение, такой клиент требуется, например, для BOINC, XTremweb и Entropia \cite{Choi2007DesktopGrid}. 



Одной из систем Desktop Grid является BOINC (Berkley Open Infrastructure for Network Computing).

Основные цели BOINC \cite{Korpela2012Seti}:

\begin{enumerate}

\item снизить порог вхождения в область публичных вычислений;

\item разделить ресурсы между проектами;

\item поддержать различные применения Desktop Grid-вычислений;

\item мотивировать участников.

\end{enumerate}

Плюсы использования BOINC:

\begin{enumerate}

\item возможность задействовать значительные вычислительные ресурсы длительное время;

\item бесплатный доступ к вычислениям;

\item хорошо подходит для постоянных исследований и переборных задач;

\item прост для участников.

\end{enumerate}

Недостатки:

\begin{enumerate}

\item нестабильная вычислительная мощность сети компьютеров;

\item различные по оборудованию компьютеры;

\item ограниченный размер одного вычислительного задания;

\item необходимость запуска дубликатов заданий, поскольку они исполняются в недоверенном окружении;

\item труднопрогнозируемый срок исполнения заданий.

\end{enumerate}



\smallskip \textit{Системы виртуализации вычислительных ресурсов} нацелены на использование виртуальных машин и контейнеров изоляции ресурсов в качестве единицы исполнения. 

Виртуализация "--- механизм абстракции физического оборудования и системных ресурсов операционной системы. 

В~настоящее время технологии виртуализации тесно связаны с концепцией облачных вычислений.  

Облачные вычисления по своей природе требуют задействовать технологии, позволяющие отделить приложения от операционной системы и реального оборудования. 

Существуют два метода программной виртуализации:

\begin{enumerate}

\item динамическая трансляция;

\item паравиртуализация.

\end{enumerate}



При динамической трансляции инструкции гостевой ОС транслируются гипервизором в безопасные инструкции, после чего управление возвращается гостевой ОС. В методе паравиртуализации требуется модифицировать исходный код ядра гостевой ОС, гипервизор предоставляет гостевой ОС специальный API для работы с системными ресурсами, такими как страницы памяти.



Помимо программной виртуализации многие гипервизоры поддерживают технологии аппаратной виртуализации, например, Intel VT-x и AMD Pacific. Аппаратная поддержка требует реализации дополнительного набора инструкций, упрощающих выполнение на аппаратном уровне операций для предоставления прямого доступа к ресурсам процессора из гостевых систем. К~системам виртуализации относятся гипервизоры: Xen, VMware и KVM. 



Приложения высокопроизводительных вычислений могут получить выигрыш от виртуализации несколькими путями \cite{Mergen2006Virtualization}: 

\begin{enumerate}

\item гипервизор может гарантировать фиксированное выделение ресурсов для гостевых ОС, таких как процессорное время и оперативная память;

\item узлы виртуального кластера могут быть запущены на различных ядрах одного или нескольких реальных узлов, причём конфигурацию узлов и их физическое расположение можно менять во время работы;

\item изоляция пользовательских процессов делает системы на базе виртуализации более устойчивыми к сбою пользовательского ПО;

\item гипервизор берёт на себя обеспечение безопасности  процессов и данных.

\end{enumerate}

Недостатки:

\begin{enumerate}

\item сложность ПО виртуализации;

\item виртуализация требует дополнительных ресурсов процессора и памяти на размещение гостевой ОС;

\item долгий (по сравнению с запуском приложения) холодный старт новых виртуальных машин.

\end{enumerate}



Стоит отдельно выделить технику контейнеризации приложений как способ изоляции приложений и частичной виртуализации ресурсов. В отличие от гипервизоров контейнерная виртуализация не требует виртуализации гостевой ОС. Она задействует механизмы ядра операционной системы, которые позволяют изолировать процессы и выделить для них физические ресурсы системы \cite{Xavier2013Performance}. Такой подход реализован в системах OpenVZ, LXC и Docker.



И хотя контейнерная виртуализация по сравнению с виртуальными машинами требует меньших ресурсов, она всё же имеет недостатки:

\begin{enumerate}

\item меньший уровень изоляции пользовательского ПО на уровне разделяемых ресурсов;

\item меньший уровень безопасности приложений и данных.

\end{enumerate}



\smallskip \textit{Платформы распределённых вычислений} представлены как отдельный \linebreak класс~систем, поскольку задействуют другие подходы к организации распределённых вычислений. Например, используют другие системы для реального исполнения задач и, самое главное, предоставляют высокий уровень абстракции от физических ресурсов.

Этому классу систем принадлежат следующие:

\begin{enumerate}

\item Globus Toolkit "--- позволяет запускать задачи с использованием планировщиков PBS, Condor, Platform LSF \cite{Foster2005Globus};

\item Hadoop "--- набор библиотек и служебных программ для организации вычислений в парадигме MapReduce \cite{Taylor2010Overview};

\item CLAVIRE "--- многопрофильная инструментально-технологическая среда, позволяющая строить распределённые приложения из готовых вычислительных узлов, определяя потоки данных \cite{Knyazkov2012Clavire}.

\end{enumerate}



CLAVIRE реализует концепцию композитных приложений, в которой наибольшее значение имеют потоки данных. Для моделирования таких потоков используется графическая нотация процессов Workflow, позволяющая определить потоки данных между частями приложения и собрать из них приложение. Моделирование потоков данных предполагает, что все необходимые вычислительные алгоритмы уже разработаны, а в облаке есть готовые (запущенные или стартующие по требованию) узлы. Поэтому такой процесс работы над приложениями непригоден для регулярного использования и исполнения массовых расчётов. Композитные приложения лучше подходят для постоянно исполняемых вычислений, например, мониторинга состояния транспортных систем или прогнозирования наводнений \cite{Ivanov2012Simulation}. 



К достоинствам платформ распределённых вычислений стоит отнести развитые средства разработки приложений, многообразие способов  запуска и мониторинг исполнения приложений.

Недостатки:

\begin{enumerate}

\item сложность отладки;

\item специализированный процесс разработки приложений, требующий задействования средств платформы (Hadoop, CLAVIRE);

\item сокрытие деталей взаимодействия приложений с оборудованием.

\end{enumerate}



\smallskip



\Section{Инструменты разработки Templet}

Веб-сервис автоматизации параллельных вычислений на суперкомпьютере <<Сергей Королёв>> Самарского государственного аэрокосмического университета  Templet состоит из 3 компонентов:

\begin{enumerate}

\item библиотеки параллельного программирования Templet SDK;

\item веб-сервис для запуска и мониторинга задач на суперкомпьютере Templet Web;

\item подсистема мониторинга состояния кластера.

\end{enumerate}



Templet SDK "--- комплект инструментов для разработки параллельных приложений. Включает в себя препроцессор исходных текстов и библиотеку времени исполнения. Препроцессор позволяет разработать и отладить последовательную программу на локальной машине, а затем запустить код на исполнение на суперкомпьютере в параллельном режиме (код будет получен автоматически из последовательной версии, для параллельной работы будет использоваться API POSIX Threads). 

Применение инструментов Templet SDK не требует от пользователя знаний методов параллельного программирования в~UNIX или Windows и позволяет сосредоточиться на решении прикладной задачи.



Веб-приложение для управления проектами и задачами Templet Web, развёрнутое по адресу \href{http://templet.ssau.ru/templet}{\tt http://templet.ssau.ru/templet}, позволяет произвести следующие действия:

\begin{itemize}

\item начать разработку проекта параллельного приложения при помощи шаблонов вычислительных методов пакета Templet SDK: портфель задач, конвейер и другие;

\item организовать частное облако окружений для запуска проектов;

\item организовать работу команды над проектом с применением VCS;

\item получать доступ к общим окружениям и запущенным задачам;

\item разворачивать приложение в тестовом и целевом окружениях;

\item отслеживать работу приложения во время продолжительных вычислений;

\item получать уведомления о статусе задач в целевых окружениях;

\item управлять бинарными зависимостями шаблонов и проектов для различных архитектур и платформ.

\end{itemize}



Шаблоны проектов позволяют предоставлять готовые структуры проектов и стандартизировать пакеты поставки приложений в целевые окружения. Проекты строятся на основе шаблонов и могут хранить код в системах контроля версий (VCS), в проект подключаются общие окружения для запуска задач. Задачи и окружения проекта доступны всей команде проекта, что позволяет эффективно организовать коллективную работу. Результаты вычислений сохраняются в архиве, к ним можно получать доступ не только по завершении задач, но и позднее. Сервис ведёт учёт использования ресурсов суперкомпьютера и предоставляет пользователям статистику по запущенным задачам и загрузке ресурсов. Это помогает оценить максимальный объём ресурсов, который сможет использовать приложение, с учётом допустимого времени ожидания вычислительной задачи в очереди пакетной системы.



Templet Web относится к классу высокоуровневых систем и задействует в качестве исполняющего механизма планировщик Torque или отдельно запускаемый процесс. В системе поддерживаются окружения под управлением ОС Linux (в том числе виртуальные) и суперкомпьютер «Сергей Королёв». Система имеет расширяемую архитектуру, что позволяет добавлять поддерживаемые окружения по мере её развития.  В Templet Web основной акцент сделан на простом процессе разработки приложений и интеграции средств разработки и платформы для исполнения. Сервис работает на двух языках "--- русском и английском, он доступен для использования не только студентам и преподавателям СГАУ, но и любому стороннему пользователю.



В основе механизма взаимодействия системы и целевых окружений лежит прямое подключение к окружению по безопасному протоколу SSH (Secure Shell). Большим преимуществом такого подхода является возможность использования неподготовленных окружений на базе ОС Linux, поскольку для работы веб-сервиса в этом случае не требуется установки дополнительного ПО и произведения настройки окружения для взаимодействия с системой. Недостатком такого решения является необходимость иметь прямое подключение от веб-сервиса к серверу окружения (или к узлу управления в случае суперкомпьютера). Зачастую высокопроизводительные кластеры недоступны для использования за пределами сети организации, выходом в таком случае может стать организация защищённой виртуальной сети (VPN). Стоит отметить, что установка дополнительного ПО (например, агента на стороне кластера) сопряжена с большими сложностями, ведь политика безопасности окружения может ограничивать возможность запуска сторонних приложений в режиме демона.



Подсистема мониторинга суперкомпьютера <<Сергей Королёв>> собирает детальную информацию по использованию суперкомпьютера:

\begin{itemize}

\item состояние очереди пакетной системы;

\item запрашиваемые задачами ресурсы;

\item состояние узлов суперкомпьютера.

\end{itemize}



Подсистема работает независимо от веб-сервиса, что позволяет производить обслуживание и обновление веб-сервиса Templet Web без остановки сбора информации о функционировании кластера. Информация, собираемая подсистемой мониторинга, используется для решения задач прогнозирования доступных вычислительных ресурсов кластера\cite{Artam2014BasesApproaches}:

\begin{itemize}

\item прогноз количества доступных вычислительных ресурсов;

\item прогноз момента запуска вычислительной задачи в очереди пакетной системы.

\end{itemize}



В основе подсистемы мониторинга используются конвейеры из подпрограмм"=агентов \cite{Agha1999Actors}. 

Каждый агент реализует небольшую хорошо инкапсулированную функциональность:

\begin{itemize}

\item получение данных о состоянии кластера;

\item разбор данных и извлечение ключевых показателей;

\item сканирование всех собранных данных;

\item создание новой записи о состоянии в базе данных (БД);

\item обновление записи о состоянии.

\end{itemize}



Такой набор микрофункциональных компонентов позволяет строить $N$ конвейеров под несколько классов задач и исполнять масштабные операции в виде $N$ параллельных конвейеров. 

Причём каждый агент исполняется конкурентно (может исполняться параллельно). Гибкая архитектура конвейеров на базе агентов используется при масштабных обновлениях БД мониторинга, например, при расчёте нового показателя для всех исторических данных. 

В~штатном режиме работы мониторинга используется один конвейер команд:

\begin{enumerate}

\item получить слепок состояния кластера;

\item выделить показатели;

\item создать новую запись в БД.

\end{enumerate}



Агенты реализованы на языке Scala \cite{Haller2009Scala}, для организации взаимодействия агентов используется библиотека Akka Actors, обеспечивающая слой абстракции и позволяющая исполнять агенты конкурентно \cite{Shams2012Actors}.



Данные подсистемы мониторинга доступны в виде графиков загрузки кластера в веб-сервисе Templet Web. Открытые данные мониторинга доступны на сайте проекта (\href{http://templet.ssau.ru/wiki/%D0%BE%D1%82%D0%BA%D1%80%D1%8B%D1%82%D1%8B%D0%B5_%D0%B4%D0%B0%D0%BD%D0%BD%D1%8B%D0%B5}{\tt http://templet.ssau.ru/wiki/открытые\_дан\-ные}).



\smallskip



\Section{Решение прикладных задач при помощи инструментов Templet}

Сервис Templet \cite{ArtamVostNazar2012IntegrationService} используется в СГАУ при работе над приложениями научных вычислений. 

Одно из приложений, в разработке которых использовался инструментарий Templet, "--- приложение расчёта параметров многомерных динамических систем. 

В~самом приложении задействован Templet SDK \cite{Vostokin2014Samstu}, запуски выполнялись через сервис Templet Web.



Задача анализа многомерных динамических систем с помощью сечений Пуанкаре содержит независимые по данным этапы вычисления координат пересечения траекторий с плоскостью сечения \cite{VostDorArtam2014GPT}. Распараллеливание вычислений одной траектории является сложной задачей, поскольку инструкции расчёта одной траектории сильно связаны по данным. Для решения этой задачи использовался шаблон Templet SDK ``Портфель задач''. Распараллеливание приложения для этого шаблона состоит в реализации функции извлечения задачи из портфеля, её обработки и помещения результатов работы в портфель. Последовательность выполнения при параллельном исполнении сохраняется только для функции извлечения задач. Эксклюзивным доступом к данным портфеля обладают функции извлечения задачи и помещения результатов в портфель.



Исходная задача была декомпозирована и представлена в виде системного (Templet SDK и шаблон ``Портфель задач'') и прикладного уровней. Шаблон вычислительного метода реализуется системными программистами и может быть повторно использован во многих проектах. Прикладные программисты или специалисты в предметной области разрабатывают приложение для вычислений в терминах и идиомах шаблона вычислительного метода, что избавляет их от необходимости ручного распараллеливания кода.



Templet SDK, включающий набор шаблонов вычислительных методов, упростил разработку параллельного приложения, которое затем запускалось при помощи веб-сервиса Templet Web. Инструментарий позволяет запускать приложение в нескольких режимах: отладка, эмуляция, приложение Win32, POSIX"=приложение, распределённое MPI"=приложение. 

Режим эмуляции позволяет оценить максимальное ускорение (оптимистичную оценку) для большого числа процессоров без реального запуска на суперкомпьютере.



\smallskip



\Section[1]{Заключение}

Инструментарий имеет большое значение при разработке приложений научной направленности, поскольку крайне важно грамотно выбрать платформу для организации разработки, средства распараллеливания и организации высокопроизводительных вычислений. Набор инструментов Templet предоставляет не только сервис для запуска приложений, но и библиотеки шаблонов параллельных приложений, средства для разработки и отладки прикладных решений. Он позволяет прикладным программистам сосредоточить усилия на работе вычислительных алгоритмов, снизить трудоёмкость разработки и сэкономить ресурсы на стадии тестирования и отладки.





%\newpage



%% Благодарности, гранты и прочее



\thanks{\textbf{Благодарности.} 

Авторы выражают благодарность Самарскому государственному аэрокосмическому университету за поддержку этого исследования. Работа выполнена при частичной поддержке Российского фонда фундаментальных исследований (проект № 15--08--05934-а) и  при государственной поддержке Министерства образования и науки РФ в рамках реализации мероприятий Программы повышения конкурентоспособности СГАУ среди ведущих мировых научно-образовательных центров на 2013--2020 годы. 

}







\thanks{{\bf ORCIDs}



{\bf Юрий Сергеевич Артамонов:}     \url{http://orcid.org/0000-0002-9301-7041}



{\bf Сергей Владимирович Востокин:} \url{http://orcid.org/0000-0001-8106-6893}



}





\RusRef



\begin{thebibliography}{99}



\Bibitem{Keyes2013Challenges}

\by Keyes~D.~E., et~al.

\paper   Multiphysics simulations: Challenges and opportunities

\yr 2013 

\jour International Journal of High Performance Computing Applications 

\vol 27 

\issue 1 

\pages 4--83 

\crossref{10.1177/1094342012468181};

%\moreref

%\preprintinfo  

Technical Report no. ANL/MCS-TM-321 Rev.~1.1,

%\yr 

2012;

\crossref{10.2172/1034263}







\Bibitem{Sullivan2009RCS}

\by Sullivan~B.

\paper Making Sense of Revision-Control Systems

\crossref{10.1145/1562164.1562183}

\jour Commun. ACM

\yr 2009

\vol 52

\issue 9

\pages 56--62





\Bibitem{Neitsch2012Build}

\by Neitsch~A., Wong~K., Godfrey~M.~W.

\crossref{10.1109/ICSM.2012.6405265}

\paper Build System Issues in   Multilanguage Software

\inbook  28th IEEE International Conference on Software Maintenance (ICSM)

\yr  2012 

\publaddr Riva del Garda, Trento, Italy

\pages 140--149





\Bibitem{Ettl2012CI}

\by Ettl~M., Neidhardt~A., Brisken~W., Dassing~R.

\paper Continuous Software   Integration and Quality Control during Software Development

\inbook Seventh General Meeting (GM2012) of the international VLBI Service for Geodesy and Astrometry (IVS)

\procinfo Madrid, Spain, March 4--9, 2012

  \eds D.~Behrend, K.~D.~Baver

\publaddr National Aeronautics and Space Administration

\yr 2012

\pages 227--230





\Bibitem{Yan2008Comparative}

\by Yan~Y., Chapman~B.

\preprint Comparative Study of Distributed Resource Management Systems – SGE, LSF, PBS Pro, and LoadLeveler

\elink \url{http://www.dcc.fc.up.pt/~ines/aulas/1213/CG/papers/RMSComparison.pdf}

\yr 2008

\totalpages 19 



\Bibitem{Anderson2004Boinc}

\by Anderson~D.~P.

\paper Boinc: A system for public-resource computing and   storage

\crossref{10.1109/grid.2004.14}

\inbook  Fifth IEEE/ACM International Workshop on Grid Computing

\yr 2004

\pages 4--10





\Bibitem{Choi2007DesktopGrid}

\by SungJin Choi, et al.

\crossref{10.1109/CCGRID.2007.31}

\paper Characterizing and  Classifying Desktop Grid

\inbook Seventh IEEE International Symposium on Cluster Computing and the Grid (CCGrid '07)

\publaddr Rio de Janeiro, Brazil

\yr 2007

\pages 743--748





\Bibitem{Korpela2012Seti}

\jour  Annual Review of Earth and Planetary Sciences

\by Korpela~E.~J.

\paper SETI@home, BOINC, and Volunteer Distributed Computing

\yr 2012 

\vol 40 

\issue 1 

\pages 69--87

\crossref{10.1146/annurev-earth-040809-152348}





\Bibitem{Mergen2006Virtualization}

\by Mergen~M.~F., Uhlig~V., Krieger~O., Xenidis~J.

\paper Virtualization for  high-performance computing

\jour  SIGOPS Oper. Syst. Rev. 

\yr 2006

\vol 40

\issue 2

\pages 8--11

\crossref{10.1145/1131322.1131328}





\Bibitem{Xavier2013Performance}

\by Xavier~M.~G.,  et~al.

\crossref{10.1109/pdp.2013.41}

\paper Performance evaluation   of container-based virtualization for high performance computing   environments 

\inbook 21st Euromicro International Conference on Parallel, Distributed, and Network-Based Processing

\yr 2013

\pages 233--240





\Bibitem{Foster2005Globus}

\by Foster~I.

\paper Globus toolkit version 4: Software for service-oriented   systems

\jour J. Comput. Sci. Technol.

\crossref{10.1007/s11390-006-0513-y}

\vol 21 

\issue 4 

\yr 2006

\pages 513--520 

\moreref

\yr 2005

\vol 3779 

\serial Lecture Notes in Computer Science 

\publ Springer 

\publaddr Berlin, Heidelberg

\crossref{10.1007/11577188_2}

\inbook  Network and Parallel Computing

\pages 2--13







\Bibitem{Taylor2010Overview}

\by Taylor~R.~C.

\paper An overview of the Hadoop, MapReduce, HBase framework  and its current applications in bioinformatics 

\jour BMC Bioinformatics

\crossref{10.1186/1471-2105-11-s12-s1}

\yr 2010

\vol 11

\issueinfo Suppl~12

\pages S1





\Bibitem{Knyazkov2012Clavire}

\by Knyazkov~K.~V., Kovalchuk~S.~V., Tchurov~T.~N.,

 Maryin~S.~V., Boukhanovsky~A.~V.

 \paper CLAVIRE:   e-Science infrastructure for data-driven computing

\jour Journal of  Computational Science

\crossref{10.1016/j.jocs.2012.08.006}

\yr 2012

\vol 3

\issue 6

\pages 504--510





\Bibitem{Ivanov2012Simulation}

\by Ivanov~S.~V., Kosukhin~S.~S., Kaluzhnaya~A.~V., Boukhanovsky~A.~V.

\paper Simulation-based collaborative decision support for surge floods prevention  in St.~Petersburg 

\jour Journal of Computational Science

\yr 2012 

\vol 3 

\issue 6 

\pages 450--455

\crossref{10.1016/j.jocs.2012.08.005}







\RBibitem{Artam2014BasesApproaches}

\by Артамонов~Ю.~С.

\paper Основные подходы прогнозирования доступных вычислительных ресурсов в кластерных системах

\inbook Перспективные информационные технологии (ПИТ 2014)

\bookinfo Труды Международной научно-технической конференции

\yr 2014

\pages 305--310

\publ Изд-во Самарского научного центра РАН

\publaddr Самара

\ed С.~А.~Прохоров

\elink \url{http://templet.ssau.ru/wiki/lib/exe/fetch.php?media=pit2014:prediction.pdf}

%\misctransl

%%\by Vostokin~S.~V.

%%\paper The basic syntax of TEMPLET markup language for representing process-per-message model

%\inbook Perspektivnye informatsionnye tekhnologii (PIT 2014) \rm [Advanced information technologies (PIT 2014)]

%\bookinfo Proceedings of the Technical Conference

%\yr 2014

%\pages 305--310

%\publ Samara Research Center

%\publaddr Samara

%\ed S.~A.~Prokhorov

%\lang In Russian





\Bibitem{Agha1999Actors}

\by Agha~G.~A., Kim~W.

\paper Actors: A unifying model for parallel and   distributed computing

\crossref{10.1016/S1383-7621(98)00067-8}

\jour Journal of   Systems Architecture

\yr 1999

\vol 45

\issue 15

\pages 1263--1277









\Bibitem{Haller2009Scala}

\by Haller~P., Odersky~M.

\paper Scala Actors: Unifying thread-based and   event-based programming

\crossref{10.1016/j.tcs.2008.09.019}

\jour Theoretical   Computer Science

\yr 2009

\vol 410 

\issue 2--3 

\pages 202--220 













%%%%









\Bibitem{Shams2012Actors}

\by Shams~M.~I., Vivek~S.

\crossref{10.1145/2384616.2384671}

\paper Integrating task parallelism  with actors

\inbook  Proceedings of the ACM international conference on Object oriented programming systems languages and applications (OOPSLA '12)

\publaddr New York, NY, USA

\yr 2012

\pages 753--772







\RBibitem{ArtamVostNazar2012IntegrationService}

\by Артамонов~Ю.~С., Востокин~С.~В., Назаров~Ю.~П.

\paper Templet – Сервис непрерывной интеграции для разработки высокопроизводительных приложений

\inbook Высокопроизводительные параллельные вычисления на кластерных системах

\bookinfo Материалы XII Всероссийской конференции

\yr 2012

\pages 82

\publ Изд-во НГУ

\publaddr Нижний Новгород

\elink \url{http://www.hpcc.unn.ru/file.php?id=713}





\RBibitem{Vostokin2014Samstu}

\by Востокин~С.~В.

\paper Препроцессор языка Templet: инструмент  программирования в терминах модели <<процесс-сообщение>>

\crossref{10.14498/vsgtu1334}

\jour Вестн. Сам. гос. техн. ун-та. Сер. Физ.-мат. науки

\issue 3(36)

\yr 2014

\pages 169--182





\RBibitem{VostDorArtam2014GPT}

\by Востокин~С.~В., Дорошин~А.~В., Артамонов~Ю.~С.

\preprint Программный комплекс   Templet. Организация прикладных вычислений на базе суперкомпьютера <<Сергей

  Королёв>>

  \yr 2014

\totalpages 5

\elink \url{http://templet.ssau.ru/wiki/_media/pit2014/templetweb.pdf}







\end{thebibliography}





\RuDate





\newpage



\makeentitle









\thanks{\textbf{Acknowledgments.} 

Our thanks go to Samara State Aerospace University (SSAU) for its continuous support to this research. 

This work was partially supported by Russian Foundation for Basic Research (project no. 15--08--05934-a) and by the Ministry of Education and Science of the Russian Federation in  the  framework  of  the  implementation  of  the  Program  of  increasing  the  competitiveness  of SSAU among the world's leading scientific and educational centers over the period from 2013 till 2020.

}



\thanks{{\bf ORCIDs}



{\bf Yuriy S. Artamonov:} \url{http://orcid.org/0000-0002-9301-7041}



{\bf Sergey V. Vostokin:} \url{http://orcid.org/0000-0001-8106-6893}



}







\EngRef







\begin{thebibliography}{99}



\Bibitem{Keyes2013Challenges:eng}

\by Keyes~D.~E., et~al.

\paper   Multiphysics simulations: Challenges and opportunities

\yr 2013 

\jour International Journal of High Performance Computing Applications 

\vol 27 

\issue 1 

\pages 4--83 

\crossref{10.1177/1094342012468181};

%\moreref

%\preprintinfo  

Technical Report no. ANL/MCS-TM-321 Rev.~1.1,

%\yr 

2012;

\crossref{10.2172/1034263}







\Bibitem{Sullivan2009RCS:eng}

\by Sullivan~B.

\paper Making Sense of Revision-Control Systems

\crossref{10.1145/1562164.1562183}

\jour Commun. ACM

\yr 2009

\vol 52

\issue 9

\pages 56--62





\Bibitem{Neitsch2012Build:eng}

\by Neitsch~A., Wong~K., Godfrey~M.~W.

\crossref{10.1109/ICSM.2012.6405265}

\paper Build System Issues in   Multilanguage Software

\inbook  28th IEEE International Conference on Software Maintenance (ICSM)

\yr  2012 

\publaddr Riva del Garda, Trento, Italy

\pages 140--149





\Bibitem{Ettl2012CI:eng}

\by Ettl~M., Neidhardt~A., Brisken~W., Dassing~R.

\paper Continuous Software   Integration and Quality Control during Software Development

\inbook Seventh General Meeting (GM2012) of the international VLBI Service for Geodesy and Astrometry (IVS)

\procinfo Madrid, Spain, March 4--9, 2012

  \eds D.~Behrend, K.~D.~Baver

\publaddr National Aeronautics and Space Administration

\yr 2012

\pages 227--230





\Bibitem{Yan2008Comparative:eng}

\by Yan~Y., Chapman~B.

\preprint Comparative Study of Distributed Resource Management Systems – SGE, LSF, PBS Pro, and LoadLeveler

\elink \url{http://www.dcc.fc.up.pt/~ines/aulas/1213/CG/papers/RMSComparison.pdf}

\yr 2008

\totalpages 19 



\Bibitem{Anderson2004Boinc:eng}

\by Anderson~D.~P.

\paper Boinc: A system for public-resource computing and   storage

\crossref{10.1109/grid.2004.14}

\inbook  Fifth IEEE/ACM International Workshop on Grid Computing

\yr 2004

\pages 4--10





\Bibitem{Choi2007DesktopGrid:eng}

\by SungJin Choi, et al.

\crossref{10.1109/CCGRID.2007.31}

\paper Characterizing and  Classifying Desktop Grid

\inbook Seventh IEEE International Symposium on Cluster Computing and the Grid (CCGrid '07)

\publaddr Rio de Janeiro, Brazil

\yr 2007

\pages 743--748





\Bibitem{Korpela2012Seti:eng}

\jour  Annual Review of Earth and Planetary Sciences

\by Korpela~E.~J.

\paper SETI@home, BOINC, and Volunteer Distributed Computing

\yr 2012 

\vol 40 

\issue 1 

\pages 69--87

\crossref{10.1146/annurev-earth-040809-152348}





\Bibitem{Mergen2006Virtualization:eng}

\by Mergen~M.~F., Uhlig~V., Krieger~O., Xenidis~J.

\paper Virtualization for  high-performance computing

\jour  SIGOPS Oper. Syst. Rev. 

\yr 2006

\vol 40

\issue 2

\pages 8--11

\crossref{10.1145/1131322.1131328}





\Bibitem{Xavier2013Performance:eng}

\by Xavier~M.~G.,  et~al.

\crossref{10.1109/pdp.2013.41}

\paper Performance evaluation   of container-based virtualization for high performance computing   environments 

\inbook 21st Euromicro International Conference on Parallel, Distributed, and Network-Based Processing

\yr 2013

\pages 233--240





\Bibitem{Foster2005Globus:eng}

\by Foster~I.

\paper Globus toolkit version 4: Software for service-oriented   systems

\jour J. Comput. Sci. Technol.

\crossref{10.1007/s11390-006-0513-y}

\vol 21 

\issue 4 

\yr 2006

\pages 513--520 

\moreref

\yr 2005

\vol 3779 

\serial Lecture Notes in Computer Science 

\publ Springer 

\publaddr Berlin, Heidelberg

\crossref{10.1007/11577188_2}

\inbook  Network and Parallel Computing

\pages 2--13







\Bibitem{Taylor2010Overview:eng}

\by Taylor~R.~C.

\paper An overview of the Hadoop, MapReduce, HBase framework  and its current applications in bioinformatics 

\jour BMC Bioinformatics

\crossref{10.1186/1471-2105-11-s12-s1}

\yr 2010

\vol 11

\issueinfo Suppl~12

\pages S1





\Bibitem{Knyazkov2012Clavire:eng}

\by Knyazkov~K.~V., Kovalchuk~S.~V., Tchurov~T.~N.,

 Maryin~S.~V., Boukhanovsky~A.~V.

 \paper CLAVIRE:   e-Science infrastructure for data-driven computing

\jour Journal of  Computational Science

\crossref{10.1016/j.jocs.2012.08.006}

\yr 2012

\vol 3

\issue 6

\pages 504--510





\Bibitem{Ivanov2012Simulation:eng}

\by Ivanov~S.~V., Kosukhin~S.~S., Kaluzhnaya~A.~V., Boukhanovsky~A.~V.

\paper Simulation-based collaborative decision support for surge floods prevention  in St.~Petersburg 

\jour Journal of Computational Science

\yr 2012 

\vol 3 

\issue 6 

\pages 450--455

\crossref{10.1016/j.jocs.2012.08.005}







\Bibitem{Artam2014BasesApproaches:eng}

\elink \url{http://templet.ssau.ru/wiki/lib/exe/fetch.php?media=pit2014:prediction.pdf}

\by Vostokin~S.~V.

\paper The basic syntax of TEMPLET markup language for representing process-per-message model

\inbook Perspektivnye informatsionnye tekhnologii (PIT 2014) \rm [Advanced information technologies (PIT 2014)]

\bookinfo Proceedings of the Technical Conference

\yr 2014

\pages 305--310

\publ Samara Research Center

\publaddr Samara

\ed S.~A.~Prokhorov

\lang In Russian





\Bibitem{Agha1999Actors:eng}

\by Agha~G.~A., Kim~W.

\paper Actors: A unifying model for parallel and   distributed computing

\crossref{10.1016/S1383-7621(98)00067-8}

\jour Journal of   Systems Architecture

\yr 1999

\vol 45

\issue 15

\pages 1263--1277









\Bibitem{Haller2009Scala:eng}

\by Haller~P., Odersky~M.

\paper Scala Actors: Unifying thread-based and   event-based programming

\crossref{10.1016/j.tcs.2008.09.019}

\jour Theoretical   Computer Science

\yr 2009

\vol 410 

\issue 2--3 

\pages 202--220 







\Bibitem{Shams2012Actors:eng}

\by Shams~M.~I., Vivek~S.

\crossref{10.1145/2384616.2384671}

\paper Integrating task parallelism  with actors

\inbook  Proceedings of the ACM international conference on Object oriented programming systems languages and applications (OOPSLA '12)

\publaddr New York, NY, USA

\yr 2012

\pages 753--772







\Bibitem{ArtamVostNazar2012IntegrationService:eng}

\elink \url{http://www.hpcc.unn.ru/file.php?id=713}

\by Artamonov~Yu.~S., Vostokin~S.~V., Nazarov~Yu.~P.

\paper Templet – Continuous integration service for the development of high performance applications

\inbook Vysokoproizvoditel'nye parallel'nye vychisleniia na klasternykh sistemakh \rm [High-performance parallel computing on cluster systems]

\bookinfo Proceedings of XII All-Russian Conference

\yr 2012

\pages 82

\publ NSU Publ.

\publaddr Nizhny Novgorod

\lang In Russian





\Bibitem{Vostokin2014Samstu:eng}

\lang In Russian

\crossref{10.14498/vsgtu1334}

\by Vostokin~S.~V.

\paper The Templet Language Preprocessor: A Programming Tool for Process-per-Message Modeling

\jour Vestn. Samar. Gos. Tekhn. Univ. Ser. Fiz.-Mat. Nauki

\yr 2014

\issue 3(36)

\pages 169--182



\Bibitem{VostDorArtam2014GPT:eng}

\lang In Russian

\by Vostokin~S.~V., Doroshin~A.~V., Artamonov~Yu.~S.

\preprint The software package Templet. Organizing applications into the supercomputer ``Sergei Korolev''

  \yr 2014

\totalpages 5

\elink \url{http://templet.ssau.ru/wiki/_media/pit2014/templetweb.pdf}







\end{thebibliography}





\EngDate



\Russian

\newpage





%\end{document}

